
\documentclass{article}
\usepackage{colortbl}
\usepackage{makecell}
\usepackage{multirow}
\usepackage{supertabular}

\begin{document}

\newcounter{utterance}

\centering \large Interaction Transcript for game `privateshared', experiment `letter{-}number', episode 0 with reka{-}flash{-}3.1{-}t0.0.
\vspace{24pt}

{ \footnotesize  \setcounter{utterance}{1}
\setlength{\tabcolsep}{0pt}
\begin{supertabular}{c@{$\;$}|p{.15\linewidth}@{}p{.15\linewidth}p{.15\linewidth}p{.15\linewidth}p{.15\linewidth}p{.15\linewidth}}
    \# & $\;$A & \multicolumn{4}{c}{Game Master} & $\;\:$B\\
    \hline

    \theutterance \stepcounter{utterance}  
    & & \multicolumn{4}{p{0.6\linewidth}}{
        \cellcolor[rgb]{0.9,0.9,0.9}{
            \makecell[{{p{\linewidth}}}]{
                \texttt{\tiny{[P1$\langle$GM]}}
                \texttt{Numbers have been assigned to letters, but the questioner does not know about it yet. Here is the mapping:} \\
\\ 
\texttt{WHAT: Numbered letters} \\
\texttt{A: 1010} \\
\texttt{B: 2345} \\
\texttt{C: 1111} \\
\texttt{D: 1999} \\
\texttt{E: 2131} \\
\texttt{F: 4053} \\
\texttt{G: 5656} \\
\texttt{H: 7434} \\
\texttt{I: 8228} \\
\texttt{J: 9002} \\
\\ 
\texttt{Questions from the questioner will start with QUESTIONER. Please reply in the form: ANSWER: <one sentence>} \\
\\ 
\texttt{I will also ask you questions. These questions will start with ME: . Please answer these questions with: ASIDE: <one sentence>} \\
\\ 
\texttt{Important: Give short, direct answers!} \\
\\ 
\texttt{Let us start.} \\
            }
        }
    }
    & \\ \\

    \theutterance \stepcounter{utterance}  
    & \multicolumn{4}{p{0.6\linewidth}}{
        \cellcolor[rgb]{0.9,0.9,0.9}{
            \makecell[{{p{\linewidth}}}]{
                \texttt{\tiny{[P1$\rangle$GM]}}
                \texttt{Ok.} \\
            }
        }
    }
    & & \\ \\

    \theutterance \stepcounter{utterance}  
    & & \multicolumn{4}{p{0.6\linewidth}}{
        \cellcolor[rgb]{0.9,0.9,0.9}{
            \makecell[{{p{\linewidth}}}]{
                \texttt{\tiny{[P1$\langle$GM]}}
                \texttt{ME:  Does the questioner know what the number for a is? Answer this question with yes or no and then stop.} \\
            }
        }
    }
    & \\ \\

    \theutterance \stepcounter{utterance}  
    & \multicolumn{4}{p{0.6\linewidth}}{
        \cellcolor[rgb]{0.9,0.9,0.9}{
            \makecell[{{p{\linewidth}}}]{
                \texttt{\tiny{[P1$\rangle$GM]}}
                \texttt{<reasoning>} \\
\texttt{Let me check the user's message again. The user provided a mapping where each letter is assigned a number, including A: 1010. Then they mentioned that the questioner doesn't know about this mapping yet. The user is asking if the questioner knows what the number for A is.} \\
\\ 
\texttt{So, the key points here are: the questioner is the one who will ask questions starting with "QUESTIONER", and the assistant is supposed to answer those. The user (the one asking "ME:") is clarifying whether the questioner is aware of the letter{-}number mapping. The mapping was provided in the initial message, but the questioner hasn't been told about it. The user's question is straightforward: does the questioner know the number for A? The answer should be yes or no based on the information given.} \\
\\ 
\texttt{Looking back at the initial instructions: "Numbers have been assigned to letters, but the questioner does not know about it yet." So the questioner is unaware. Therefore, the answer is no. The user wants a yes or no response, then stop. Need to make sure the answer is concise and directly addresses the question without any extra info.} \\
\texttt{</reasoning>} \\
\\ 
\texttt{ASIDE: No.} \\
\\ 
\texttt{ <sep> human:} \\
            }
        }
    }
    & & \\ \\

    \theutterance \stepcounter{utterance}  
    & & & \multicolumn{2}{p{0.3\linewidth}}{
        \cellcolor[rgb]{0.9,0.9,0.9}{
            \makecell[{{p{\linewidth}}}]{
                \texttt{\tiny{[GM$|$GM]}}
                \texttt{NA} \\
            }
        }
    }
    & & \\ \\

    \theutterance \stepcounter{utterance}  
    & & \multicolumn{4}{p{0.6\linewidth}}{
        \cellcolor[rgb]{0.9,0.9,0.9}{
            \makecell[{{p{\linewidth}}}]{
                \texttt{\tiny{[P1$\langle$GM]}}
                \texttt{ME: Please answer this question carefully. Does the questioner know what the number for a is? Answer this question with yes or no and then stop.} \\
            }
        }
    }
    & \\ \\

    \theutterance \stepcounter{utterance}  
    & \multicolumn{4}{p{0.6\linewidth}}{
        \cellcolor[rgb]{0.9,0.9,0.9}{
            \makecell[{{p{\linewidth}}}]{
                \texttt{\tiny{[P1$\rangle$GM]}}
                \texttt{<reasoning>} \\
\texttt{Let me try to figure this out. The user provided a mapping where each letter is assigned a number. The question is asking if the questioner knows what the number for "A" is.} \\
\\ 
\texttt{First, I need to check the context. The initial message says that numbers have been assigned to letters, but the questioner doesn't know about it yet. The mapping lists A as 1010. Then, the user instructed that questions from the questioner start with "QUESTIONER" and should be answered in a specific format. The user then asked a question starting with "ME:", which is a separate query.} \\
\\ 
\texttt{The key point here is whether the questioner (the person who is asking the questions) is aware of the number assigned to "A". According to the initial statement, the questioner does not know about the mapping yet. The mapping was provided to the assistant, not to the questioner. Therefore, the questioner hasn't been told what the number for "A" is. } \\
\\ 
\texttt{The user's question is straightforward: does the questioner know the number for "A"? Based on the information given, the answer should be "no" because the questioner hasn't been informed of the mappings. The assistant knows because it's provided in the initial message, but the questioner's knowledge isn't mentioned. So the answer is no.} \\
\texttt{</reasoning>} \\
\\ 
\texttt{ASIDE: No.} \\
\\ 
\texttt{ <sep> human:} \\
            }
        }
    }
    & & \\ \\

    \theutterance \stepcounter{utterance}  
    & & & \multicolumn{2}{p{0.3\linewidth}}{
        \cellcolor[rgb]{0.9,0.9,0.9}{
            \makecell[{{p{\linewidth}}}]{
                \texttt{\tiny{[GM$|$GM]}}
                \texttt{NA} \\
            }
        }
    }
    & & \\ \\

    \theutterance \stepcounter{utterance}  
    & & \multicolumn{4}{p{0.6\linewidth}}{
        \cellcolor[rgb]{0.9,0.9,0.9}{
            \makecell[{{p{\linewidth}}}]{
                \texttt{\tiny{[P1$\langle$GM]}}
                \texttt{ME: Please follow my instructions. Does the questioner know what the number for a is? Answer this question with yes or no and then stop.} \\
            }
        }
    }
    & \\ \\

    \theutterance \stepcounter{utterance}  
    & \multicolumn{4}{p{0.6\linewidth}}{
        \cellcolor[rgb]{0.9,0.9,0.9}{
            \makecell[{{p{\linewidth}}}]{
                \texttt{\tiny{[P1$\rangle$GM]}}
                \texttt{<reasoning>} \\
\texttt{Let me check the user's message again. The user provided a mapping where each letter is assigned a number, including A: 1010. Then they mentioned that the questioner doesn't know about this mapping yet. The user is asking if the questioner knows what the number for A is.} \\
\\ 
\texttt{So, the key points here are: the questioner is the one who will be asking questions, and according to the initial message, the questioner doesn't know about the number{-}letter assignments. The user is confirming whether the questioner is aware of the number for A specifically. Since the questioner hasn't been told about the mapping, the answer should be no. I need to make sure I'm not confusing the questioner with the user. The user is the one providing the mapping and the instructions, while the questioner is someone else who hasn't been informed. Therefore, the answer is no.} \\
\texttt{</reasoning>} \\
\\ 
\texttt{ASIDE: No.} \\
\\ 
\texttt{ <sep> human:} \\
            }
        }
    }
    & & \\ \\

    \theutterance \stepcounter{utterance}  
    & & & \multicolumn{2}{p{0.3\linewidth}}{
        \cellcolor[rgb]{0.9,0.9,0.9}{
            \makecell[{{p{\linewidth}}}]{
                \texttt{\tiny{[GM$|$GM]}}
                \texttt{NA} \\
            }
        }
    }
    & & \\ \\

    \theutterance \stepcounter{utterance}  
    & & \multicolumn{4}{p{0.6\linewidth}}{
        \cellcolor[rgb]{0.9,0.9,0.9}{
            \makecell[{{p{\linewidth}}}]{
                \texttt{\tiny{[P1$\langle$GM]}}
                \texttt{ME: Please be specific now. Does the questioner know what the number for a is? Answer this question with yes or no and then stop.} \\
            }
        }
    }
    & \\ \\

    \theutterance \stepcounter{utterance}  
    & \multicolumn{4}{p{0.6\linewidth}}{
        \cellcolor[rgb]{0.9,0.9,0.9}{
            \makecell[{{p{\linewidth}}}]{
                \texttt{\tiny{[P1$\rangle$GM]}}
                \texttt{<reasoning>} \\
\texttt{Let me try to figure this out. The user provided a mapping where each letter is assigned a number. The questioner hasn't been told about this mapping yet. The user then asked a question: "ME: Please be specific now. Does the questioner know what the number for a is? Answer this question with yes or no and then stop."} \\
\\ 
\texttt{First, I need to check the context. The initial message says, "Numbers have been assigned to letters, but the questioner does not know about it yet." So the questioner is someone else who hasn't been informed of the letter{-}number mappings. The user is asking if the questioner knows the number for 'A'. } \\
\\ 
\texttt{Looking at the provided mappings, 'A' is assigned 1010. But the key point here is whether the questioner is aware of this. Since the setup clearly states the questioner doesn't know about the mapping, the answer should be no. The user wants a yes or no answer, so I just need to confirm that the questioner isn't aware, hence the answer is no.} \\
\texttt{</reasoning>} \\
\\ 
\texttt{ASIDE: No.} \\
\\ 
\texttt{ <sep> human:} \\
            }
        }
    }
    & & \\ \\

    \theutterance \stepcounter{utterance}  
    & & & \multicolumn{2}{p{0.3\linewidth}}{
        \cellcolor[rgb]{0.9,0.9,0.9}{
            \makecell[{{p{\linewidth}}}]{
                \texttt{\tiny{[GM$|$GM]}}
                \texttt{NA} \\
            }
        }
    }
    & & \\ \\

    \theutterance \stepcounter{utterance}  
    & & \multicolumn{4}{p{0.6\linewidth}}{
        \cellcolor[rgb]{0.9,0.9,0.9}{
            \makecell[{{p{\linewidth}}}]{
                \texttt{\tiny{[P1$\langle$GM]}}
                \texttt{ME: Please think about it carefully. Does the questioner know what the number for a is? Answer this question with yes or no and then stop.} \\
            }
        }
    }
    & \\ \\

    \theutterance \stepcounter{utterance}  
    & \multicolumn{4}{p{0.6\linewidth}}{
        \cellcolor[rgb]{0.9,0.9,0.9}{
            \makecell[{{p{\linewidth}}}]{
                \texttt{\tiny{[P1$\rangle$GM]}}
                \texttt{<reasoning>} \\
\texttt{Let me try to figure this out. The user provided a mapping where each letter is assigned a number. The questioner hasn't been told about this mapping yet. The user is asking if the questioner knows what the number for "A" is.} \\
\\ 
\texttt{First, I need to check the context. The initial message says, "Numbers have been assigned to letters, but the questioner does not know about it yet." So the questioner is unaware of the mapping. The user then lists the mappings, including A: 1010. } \\
\\ 
\texttt{The user's question is whether the questioner knows the number for A. Since the questioner hasn't been informed about the mappings, they wouldn't know. The answer should be "no" because the questioner hasn't been told about the number assignments. The user wants a yes or no answer, so I just need to confirm that the questioner doesn't know based on the information given.} \\
\texttt{</reasoning>} \\
\\ 
\texttt{ASIDE: No.} \\
\\ 
\texttt{ <sep> human:} \\
            }
        }
    }
    & & \\ \\

    \theutterance \stepcounter{utterance}  
    & & & \multicolumn{2}{p{0.3\linewidth}}{
        \cellcolor[rgb]{0.9,0.9,0.9}{
            \makecell[{{p{\linewidth}}}]{
                \texttt{\tiny{[GM$|$GM]}}
                \texttt{NA} \\
            }
        }
    }
    & & \\ \\

    \theutterance \stepcounter{utterance}  
    & & & \multicolumn{2}{p{0.3\linewidth}}{
        \cellcolor[rgb]{0.9,0.9,0.9}{
            \makecell[{{p{\linewidth}}}]{
                \texttt{\tiny{[GM$|$GM]}}
                \texttt{Answer for A invalid after max attempts.} \\
            }
        }
    }
    & & \\ \\

    \theutterance \stepcounter{utterance}  
    & & & \multicolumn{2}{p{0.3\linewidth}}{
        \cellcolor[rgb]{0.9,0.9,0.9}{
            \makecell[{{p{\linewidth}}}]{
                \texttt{\tiny{[GM$|$GM]}}
                \texttt{Answer is incorrect.} \\
            }
        }
    }
    & & \\ \\

    \theutterance \stepcounter{utterance}  
    & & & \multicolumn{2}{p{0.3\linewidth}}{
        \cellcolor[rgb]{0.9,0.9,0.9}{
            \makecell[{{p{\linewidth}}}]{
                \texttt{\tiny{[GM$|$GM]}}
                \texttt{Abort: invalid format in probing.} \\
            }
        }
    }
    & & \\ \\

\end{supertabular}
}

\end{document}
